% ENGLISH TEMPLATE: BIBLIOGRAPHY FIX
% PREAMBLE
% Compile main.tex with XeLaTeX.
% Keep one option active in each settings group.

% LANGUAGE
% English (default):
\usepackage[brazilian,english]{babel}
% Brazilian Portuguese (replace the line above):
% \usepackage[english,brazilian]{babel}
%
% The last language is the main language.
% Babel sets hyphenation and automatic labels.
% Translate your own headings and text manually.
% For a short passage in another language:
% \foreignlanguage{brazilian}{Portuguese text}
% For a longer passage:
% \begin{otherlanguage}{brazilian}
% [Your Portuguese text here.]
% \end{otherlanguage}

% FONT
\usepackage{fontspec}
\setmainfont{Linux Libertine O}
% Use XeLaTeX, not pdfLaTeX.
% Do not load inputenc or fontenc here.

% MARGINS
% These are template design choices.
\usepackage[
  a4paper,
  margin=2.5cm,
  headheight=15pt,
  footskip=1.2cm
]{geometry}

% LINE SPACING
\usepackage{setspace}
% FAPESP scholarship projects: double spacing.
\doublespacing
% General use: replace the line above with:
% \onehalfspacing
% Or, for single spacing:
% \singlespacing
%
% Check the rules for your specific funding call.
% This setting is not universal FAPESP compliance.
% General scholarship guidance:
% https://fapesp.br/253/projeto-de-pesquisa
% Body size: 12 pt, set in main.tex.
% Cover pages use local single spacing.

% PARAGRAPHS
% Indented, with no extra space between paragraphs:
\usepackage{indentfirst}
\setlength{\parindent}{1.25cm}
\setlength{\parskip}{0pt}
% Alternative: replace the two settings above.
% \setlength{\parindent}{0pt}
% \setlength{\parskip}{0.6em}

% GRAPHICS AND COVER COLORS
\usepackage{graphicx}
\graphicspath{{logos/}{fig/}}
\usepackage[table,svgnames]{xcolor}
\usepackage{tikz}
\usetikzlibrary{
  calc,intersections,decorations.text
}
\definecolor{feecDarkBlue}{HTML}{00427A}
\definecolor{feecNavy}{HTML}{102B4E}
\definecolor{feecPurple}{HTML}{562877}
\definecolor{feecGreen}{HTML}{B0D361}
\definecolor{feecLavender}{HTML}{EACCFF}
% Color names are inherited from the original cover.

% COVER B COLORS
% Original attachment omitted these definitions.
% These editable defaults match the shared palette.
\definecolor{aimsBlue}{HTML}{00427A}
\definecolor{aimsYellow}{HTML}{F2C94C}
\definecolor{aimsTeal}{HTML}{238A8D}

% FUTURE COVERS
% Add new named colors and TikZ libraries here.
% Put each design in its own titlepage_*.tex file.
% Do not change global fonts/spacing for one cover.

% TABLES, FIGURES AND MATHEMATICS
\usepackage{array}
\usepackage{tabularx}
\usepackage{booktabs}
\usepackage{longtable}
\usepackage{float}
\usepackage{caption}
\usepackage{subcaption}
\usepackage{amsmath,amssymb}
\usepackage{enumitem}
\captionsetup{font=small,labelfont=bf}
\renewcommand{\arraystretch}{1.2}
\setcounter{secnumdepth}{3}
\setcounter{tocdepth}{2}

% LINKS
% Load BEFORE abntex2cite.
% abntex2cite detects hyperref when it loads.
\usepackage{hyperref}
\hypersetup{
  colorlinks=true,
  linkcolor=feecDarkBlue,
  citecolor=feecDarkBlue,
  urlcolor=feecDarkBlue
}

% ABNT AUTHOR-DATE CITATIONS
% Keep this package for either document language.
% It provides native \citeonline support.
\usepackage[alf]{abntex2cite}
% Do not also load cite, natbib or biblatex.
% The package selects abntex2-alf automatically.
% Do not add another \bibliographystyle command.
%
% Replace 'key' with a real entry from your .bib:
% According to \citeonline{key}, ...
% ... \cite{key}.
% \citeonline[p.~25]{key}
% \cite[p.~25]{key}
% \cite{key1,key2}
%
% Enable the bibliography at the end of main.tex.
% Use BibTeX, not Biber.
% Local compilation sequence:
% 1. xelatex main
% 2. bibtex main
% 3. xelatex main
% 4. xelatex main
% Overleaf runs these steps automatically.
%
% After switching from biblatex/Biber, use:
% Recompile from scratch.
% This clears incompatible .aux and .bbl files.
% The package uses the traditional abnTeX2 style.
% Later ABNT revisions are not applied automatically.

% FOOTER
% Language-neutral page numbering.
\usepackage{fancyhdr}
\usepackage{lastpage}
\pagestyle{fancy}
\fancyhf{}
\renewcommand{\headrulewidth}{0pt}
\renewcommand{\footrulewidth}{0.4pt}
\fancyfoot[R]{%
  \small\color{black!70}%
  \thepage\ / \pageref{LastPage}}
\fancypagestyle{plain}{%
  \fancyhf{}
  \renewcommand{\headrulewidth}{0pt}
  \renewcommand{\footrulewidth}{0.4pt}
  \fancyfoot[R]{%
    \small\color{black!70}%
    \thepage\ / \pageref{LastPage}}}
